---
tipo: Discursiva
dificuldade: Fácil
disciplina: Matemática
assunto: Relações Métricas no Triângulo Retângulo
tags: [mediana, hipotenusa, catetos]
fonte: material didático
gabarito:
resposta: A mediana mede 5 m
---
\question
Calcule a mediana relativa à hipotenusa de um triângulo cujos catetos medem 6 m e 8 m.

---
tipo: Discursiva
dificuldade: Média
disciplina: Matemática
assunto: Relações Métricas no Triângulo Retângulo
tags: [bissetriz, ângulo reto, catetos, hipotenusa]
fonte: material didático
gabarito:
resposta: Os catetos medem 9 m e 12 m
---
\question
Num triângulo retângulo, a bissetriz do ângulo reto determina sobre o lado oposto segmentos proporcionais a 3 e 4. Sabendo que a hipotenusa mede 15 m, calcule os catetos.

---
tipo: Discursiva
dificuldade: Média
disciplina: Matemática
assunto: Relações Métricas no Triângulo Retângulo
tags: [altura, catetos, hipotenusa, figure]
fonte: material didático
gabarito:
resposta:
---
\question
Determine $x$ nos casos:
\begin{parts}
\part
% figura presente no enunciado
\part
% figura presente no enunciado
\part
% figura presente no enunciado
\part
% figura presente no enunciado
\end{parts}

---
tipo: Discursiva
dificuldade: Média
disciplina: Matemática
assunto: Relações Métricas no Triângulo Retângulo
tags: [altura, catetos, projeções]
fonte: material didático
gabarito:
resposta:
---
\question
Na figura abaixo as medidas estão em centímetros. Determine $x$, $y$, $z$ e $t$.

% figura presente no enunciado

---
tipo: Discursiva
dificuldade: Média
disciplina: Matemática
assunto: Relações Métricas no Triângulo Retângulo
tags: [altura, catetos, projeções]
fonte: material didático
gabarito:
resposta:
---
\question
Determine $x$ e $y$ nos casos abaixo:
\begin{parts}
\part
% figura presente no enunciado
\part
% figura presente no enunciado
\end{parts}

---
tipo: Discursiva
dificuldade: Média
disciplina: Matemática
assunto: Relações Métricas no Triângulo Retângulo
tags: [altura, catetos, projeções]
fonte: material didático
gabarito:
resposta:
---
\question
Determine $x$ nos casos abaixo:
\begin{parts}
\part
% figura presente no enunciado
\part
% figura presente no enunciado
\end{parts}

---
tipo: Múltipla Escolha
dificuldade: Média
disciplina: Matemática
assunto: Relações Métricas no Triângulo Retângulo
tags: [hexágonos regulares, distância, velocidade, UERJ]
fonte: concurso
banca: UERJ
ano: 2009
gabarito: D
---
\question
Um piso plano é revestido de hexágonos regulares congruentes cujo lado mede 10 cm. Na ilustração de parte desse piso, T, M e F são vértices comuns a três hexágonos e representam os pontos nos quais se encontram, respectivamente, um torrão de açúcar, uma mosca e uma formiga.

% figura presente no enunciado

Ao perceber o açúcar, os dois insetos partem no mesmo instante, com velocidades constantes, para alcançá-lo. Admita que a mosca leve 10 segundos para atingir o ponto T. Despreze o espaçamento entre os hexágonos e as dimensões dos animais.
A menor velocidade, em centímetros por segundo, necessária para que a formiga chegue ao ponto T no mesmo instante em que a mosca, é igual a:
\begin{choices}
\choice 3,5
\choice 5,0
\choice 5,5
\correctchoice 7,0
\end{choices}

---
tipo: Discursiva
dificuldade: Média
disciplina: Matemática
assunto: Lei dos Senos
tags: [bissetriz, triângulo, ângulos, CEFET]
fonte: concurso
banca: CEFET
ano: 2009
gabarito:
resposta:
---
\question
Num triângulo $ABC$, sabe-se que o ângulo $\hat{A} = 60°$, que o lado $AB$ mede $2$ e a bissetriz do ângulo $\hat{A}$ tem medida $2$. Determine:
\begin{parts}
\part Os outros lados do triângulo $ABC$.
\part Os outros ângulos do triângulo $ABC$.
\end{parts}

---
tipo: Discursiva
dificuldade: Média
disciplina: Matemática
assunto: Lei dos Cossenos
tags: [natureza do triângulo, Clairaut, cosseno]
fonte: material didático
gabarito:
resposta:
---
\question
Os lados de um triângulo medem 3, 4 e 6. Indique a natureza do triângulo (acutângulo, retângulo ou obtusângulo) e calcule o cosseno do seu maior ângulo interno.

---
tipo: Múltipla Escolha
dificuldade: Média
disciplina: Matemática
assunto: Lei dos Senos
tags: [circunferência circunscrita, raio, seno]
fonte: material didático
gabarito: E
---
\question
O triângulo $ABC$ está inscrito em um círculo de raio $R$. Se $\cos\hat{A} = \dfrac{3}{5}$, o comprimento do lado $BC$ é:
\begin{choices}
\choice $\dfrac{2R}{5}$
\choice $\dfrac{3R}{5}$
\choice $\dfrac{4R}{5}$
\choice $\dfrac{6R}{5}$
\correctchoice $\dfrac{8R}{5}$
\end{choices}

---
tipo: Múltipla Escolha
dificuldade: Média
disciplina: Matemática
assunto: Relações Métricas no Triângulo Retângulo
tags: [lampião, cordas perpendiculares, distância, UFRS]
fonte: concurso
banca: UFRS
gabarito: E
---
\question
O lampião representado na figura está suspenso por duas cordas perpendiculares presas ao teto.

% figura presente no enunciado

Sabendo que essas cordas medem $\dfrac{1}{2}$ e $\dfrac{6}{5}$, a distância do lampião ao teto é:
\begin{choices}
\choice 1,69
\choice 1,3
\choice 0,6
\choice $\dfrac{1}{2}$
\correctchoice $\dfrac{6}{13}$
\end{choices}

---
tipo: Discursiva
dificuldade: Média
disciplina: Matemática
assunto: Teorema de Pitágoras
tags: [segmentos perpendiculares, distância, UFF]
fonte: concurso
banca: UFF
gabarito:
resposta:
---
\question
Os segmentos de reta $AB$, $BC$, $CD$ e $DE$ são tais que $AB \perp BC$, $BC \perp CD$ e $CD \perp DE$. As medidas de $AB$, $BC$, $CD$ e $DE$ são, respectivamente, 3 m, 4 m, 1 m e 4 m. Determine a medida do segmento $AE$.

---
tipo: Discursiva
dificuldade: Média
disciplina: Matemática
assunto: Lei dos Cossenos
tags: [circunferência, pontos, distância, UFF]
fonte: concurso
banca: UFF
gabarito:
resposta:
---
\question
Considere os pontos $P$ e $Q$ pertencentes à circunferência de centro na origem e raio 1, conforme a figura.

% figura presente no enunciado

Determine a distância entre $P$ e $Q$.

---
tipo: Discursiva
dificuldade: Média
disciplina: Matemática
assunto: Relações Métricas no Triângulo Retângulo
tags: [triângulo retângulo, projeções, altura]
fonte: material didático
gabarito:
resposta:
---
\question
Calcule o valor de $y$ na figura abaixo.

% figura presente no enunciado

---
tipo: Discursiva
dificuldade: Fácil
disciplina: Matemática
assunto: Teorema de Pitágoras
tags: [velocidade, distância, ciclistas]
fonte: material didático
gabarito:
resposta:
---
\question
Dois ciclistas partem de uma mesma cidade em direção reta: um em direção Leste e outro em direção Norte. Determine a distância que os separa duas horas após, sabendo que as velocidades dos ciclistas são de 30 km/h e 45 km/h, respectivamente.

---
tipo: Discursiva
dificuldade: Média
disciplina: Matemática
assunto: Trigonometria
tags: [ângulo de elevação, edifício, tangente]
fonte: material didático
gabarito:
resposta:
---
\question
Um observador vê um edifício, construído em terreno plano, sob um ângulo de $60°$. Se ele se afastar do edifício mais 30 m, passará a vê-lo sob ângulo de $45°$. Calcule a altura do edifício.

---
tipo: Discursiva
dificuldade: Média
disciplina: Matemática
assunto: Lei dos Senos
tags: [ponte, rio, ângulos, UFPE]
fonte: concurso
banca: UFPE
ano: 2004
gabarito:
resposta:
---
\question
Uma ponte deve ser construída sobre um rio, unindo os pontos $A$ e $B$. Para calcular o comprimento $AB$, escolhe-se um ponto $C$, na mesma margem em que $B$ está, e medem-se os ângulos $\widehat{CBA} = 57°$ e $\widehat{ACB} = 59°$.

% figura presente no enunciado

Sabendo que $BC$ mede 30 m, indique, em metros, a distância $AB$. (Dado: use as aproximações $\operatorname{sen}(59°) = 0{,}87$ e $\operatorname{sen}(64°) = 0{,}90$)

---
tipo: Discursiva
dificuldade: Difícil
disciplina: Matemática
assunto: Lei dos Senos
tags: [circunferência, raio, segmento, UNICAMP]
fonte: concurso
banca: UNICAMP
ano: 2005
gabarito:
resposta:
---
\question
Sejam $A$, $B$, $C$ e $N$ quatro pontos em um mesmo plano, conforme a figura.

% figura presente no enunciado

\begin{parts}
\part Calcule o raio da circunferência que passa pelos pontos $A$, $B$ e $N$.
\part Calcule o comprimento do segmento $NB$.
\end{parts}

---
tipo: Discursiva
dificuldade: Difícil
disciplina: Matemática
assunto: Lei dos Cossenos
tags: [demonstração, lei dos cossenos, UFRJ]
fonte: concurso
banca: UFRJ
ano: 2002
gabarito:
resposta:
---
\question
Considerando o triângulo da figura a seguir, demonstre que $a^2 = b^2 + c^2 - 2bc\cos\theta$.

% figura presente no enunciado

---
tipo: Discursiva
dificuldade: Difícil
disciplina: Matemática
assunto: Lei dos Cossenos
tags: [bissetriz, altura, FUVEST]
fonte: concurso
banca: FUVEST
ano: 2004
gabarito:
resposta:
---
\question
Um triângulo $ABC$ tem lados de comprimentos $AB = 5$, $BC = 4$ e $AC = 2$. Sejam $M$ e $N$ os pontos de $AB$ tais que $CM$ é a bissetriz relativa ao ângulo $\widehat{ACB}$ e $CN$ é a altura relativa ao lado $AB$. Determine o comprimento de $MN$.

---
tipo: Múltipla Escolha
dificuldade: Difícil
disciplina: Matemática
assunto: Relações Métricas no Triângulo Retângulo
tags: [triângulo isósceles, quadrado inscrito, FUVEST]
fonte: concurso
banca: FUVEST
gabarito: B
---
\question
Na figura abaixo, $ABC$ é um triângulo isósceles e retângulo em $A$ e $PQRS$ é um quadrado de lado $\dfrac{2\sqrt{2}}{3}$.

% figura presente no enunciado

Então, a medida do lado $AB$ é:
\begin{choices}
\choice 1
\correctchoice 2
\choice 3
\choice 4
\choice 5
\end{choices}

---
tipo: Discursiva
dificuldade: Fácil
disciplina: Matemática
assunto: Relações Métricas no Triângulo Retângulo
tags: [altura, hipotenusa, catetos]
fonte: material didático
gabarito:
resposta:
---
\question
Em um triângulo retângulo, os catetos medem 9 cm e 12 cm. Calcule a altura relativa à hipotenusa.

---
tipo: Discursiva
dificuldade: Média
disciplina: Matemática
assunto: Relações Métricas no Triângulo Retângulo
tags: [perímetro, triângulo retângulo, altitude]
fonte: material didático
gabarito:
resposta:
---
\question
Encontre o perímetro do triângulo $ABC$ abaixo, sabendo que $AH = 2$ m e $CH = 4$ m.

% figura presente no enunciado

---
tipo: Múltipla Escolha
dificuldade: Difícil
disciplina: Matemática
assunto: Trigonometria
tags: [bissetriz, incentro, triângulo retângulo, ITA]
fonte: concurso
banca: ITA
gabarito: B
---
\question
Num triângulo $ABC$, retângulo em $A$, temos $\hat{B} = 60°$. As bissetrizes dos ângulos $\hat{A}$ e $\hat{B}$ se encontram num ponto $D$. Se o segmento $BD$ mede 1 cm, então a hipotenusa mede:
\begin{choices}
\choice $\dfrac{1 + \sqrt{3}}{2}$ cm
\correctchoice $1 + \sqrt{3}$ cm
\choice $2 + \sqrt{3}$ cm
\choice $1 + 2\sqrt{2}$ cm
\choice n.d.a.
\end{choices}

---
tipo: Múltipla Escolha
dificuldade: Média
disciplina: Matemática
assunto: Lei dos Cossenos
tags: [livro, diagonais, ângulo, FUVEST]
fonte: concurso
banca: FUVEST
ano: 2002
gabarito: B
---
\question
As páginas de um livro medem 1 dm de base e $\left(1 + \sqrt{3}\right)$ dm de altura. Se este livro foi parcialmente aberto de tal forma que o ângulo entre duas páginas seja $60°$, a medida do ângulo $\alpha$, formado pelas diagonais das páginas, será:
\begin{choices}
\choice $15°$
\correctchoice $30°$
\choice $45°$
\choice $60°$
\choice $75°$
\end{choices}

---
tipo: Discursiva
dificuldade: Difícil
disciplina: Matemática
assunto: Lei dos Senos
tags: [circunferência, triângulo inscrito, PUC-Rio]
fonte: concurso
banca: PUC-Rio
ano: 2011
gabarito:
resposta:
---
\question
Considere o triângulo $ABC$ inscrito na circunferência de raio 1 com ângulos $\widehat{BAC} = 60°$ e $\widehat{ABC} = 45°$, conforme a figura abaixo.

% figura presente no enunciado

\begin{parts}
\part Calcule o comprimento de cada um dos três lados do triângulo $ABC$.
\part Calcule a área do triângulo $ABC$.
\end{parts}

---
tipo: Múltipla Escolha
dificuldade: Média
disciplina: Matemática
assunto: Trigonometria
tags: [trabalho, força, gráfico, UERJ]
fonte: concurso
banca: UERJ
ano: 2012
gabarito: D
---
\question
Uma pessoa empurrou um carro por uma distância de 26 m, aplicando uma força $F$ de mesma direção e sentido do deslocamento desse carro. O gráfico abaixo representa a variação da intensidade de $F$, em newtons, em função do deslocamento $d$, em metros.

% figura presente no enunciado

Desprezando o atrito, o trabalho total, em joules, realizado por $F$, equivale a:
\begin{choices}
\choice 117
\choice 130
\choice 143
\correctchoice 156
\end{choices}

---
tipo: Discursiva
dificuldade: Média
disciplina: Matemática
assunto: Relações Métricas no Triângulo Retângulo
tags: [pipa, quadrilátero, diagonais perpendiculares, UERJ]
fonte: concurso
banca: UERJ
ano: 2012
gabarito:
resposta:
---
\question
Para construir a pipa representada na figura abaixo pelo quadrilátero $ABCD$, foram utilizadas duas varetas, linha e papel.

% figura presente no enunciado

As varetas estão representadas pelos segmentos $AC$ e $BD$. A linha utilizada liga as extremidades $A$, $B$, $C$ e $D$ das varetas, e o papel reveste a área total da pipa. Os segmentos $AC$ e $BD$ são perpendiculares em $E$, e os ângulos $\widehat{ABC}$ e $\widehat{ADC}$ são retos. Se os segmentos $AE$ e $EC$ medem, respectivamente, 18 cm e 32 cm, determine o comprimento total da linha, representado por $AB + BC + CD + DA$.

---
tipo: Discursiva
dificuldade: Média
disciplina: Matemática
assunto: Relações Métricas no Triângulo Retângulo
tags: [triângulo retângulo, catetos, altura, projeções]
fonte: material didático
gabarito:
resposta:
---
\question
No triângulo retângulo abaixo, descubra as medidas de $m$, $n$, $b$ e $h$.

% figura presente no enunciado

---
tipo: Discursiva
dificuldade: Média
disciplina: Matemática
assunto: Lei dos Cossenos
tags: [ponteiros, relógio, UFRJ]
fonte: concurso
banca: UFRJ
ano: 2001
gabarito:
resposta:
---
\question
Os ponteiros de um relógio circular medem, do centro às extremidades, 2 metros o dos minutos e 1 metro o das horas. Determine a distância entre as extremidades dos ponteiros quando o relógio marca 4 horas.

---
tipo: Discursiva
dificuldade: Média
disciplina: Matemática
assunto: Lei dos Cossenos
tags: [triângulo, segmentos, cosseno]
fonte: material didático
gabarito:
resposta:
---
\question
Na figura abaixo, tem-se $AC = 3$, $AB = 4$ e $BC = 6$.

% figura presente no enunciado

O valor de $CD$ é:

---
tipo: Discursiva
dificuldade: Média
disciplina: Matemática
assunto: Lei dos Senos
tags: [triângulo, seno, ângulo]
fonte: material didático
gabarito:
resposta:
---
\question
No triângulo $ABC$ da figura abaixo, as medidas estão em centímetros.

% figura presente no enunciado

Calcule o seno do ângulo $\alpha$.

---
tipo: Discursiva
dificuldade: Difícil
disciplina: Matemática
assunto: Teorema de Pitágoras
tags: [triângulos equiláteros, distância, UFF]
fonte: concurso
banca: UFF
gabarito:
resposta:
---
\question
Na figura abaixo, os triângulos $ABC$ e $DEF$ são equiláteros. Sabendo que $AB$, $CD$ e $DE$ medem, respectivamente, 6 m, 4 m e 4 m, calcule a medida de $BE$.

% figura presente no enunciado

---
tipo: Discursiva
dificuldade: Difícil
disciplina: Matemática
assunto: Relações Métricas no Triângulo Retângulo
tags: [triângulo retângulo, ponto médio, perpendicular, UFF]
fonte: concurso
banca: UFF
gabarito:
resposta:
---
\question
O triângulo $PQR$ é retângulo em $Q$, $N$ é ponto médio de $QR$ e $NM$ é perpendicular a $PR$, conforme a figura abaixo.

% figura presente no enunciado

Determine a medida de $NM$.

---
tipo: Discursiva
dificuldade: Difícil
disciplina: Matemática
assunto: Teorema de Pitágoras
tags: [triângulo equilátero, quadrado, distância, UFRJ]
fonte: concurso
banca: UFRJ
gabarito:
resposta:
---
\question
Na figura, o triângulo $AEC$ é equilátero e $ABCD$ é um quadrado de lado 2 cm. Calcule a distância $BE$.

% figura presente no enunciado

---
tipo: Discursiva
dificuldade: Difícil
disciplina: Matemática
assunto: Relações Métricas no Triângulo Retângulo
tags: [quadrado, ponto médio, retângulo, UERJ]
fonte: concurso
banca: UERJ
ano: 2010
gabarito:
resposta:
---
\question
Observe a figura abaixo, que representa um quadrado $ABCD$, de papel, no qual $M$ e $N$ são os pontos médios de dois de seus lados. Esse quadrado foi dividido em quatro partes para formar um jogo.

% figura presente no enunciado

O jogo consiste em montar, com todas essas partes, um retângulo cuja base seja maior que a altura. O retângulo $PQRS$, mostrado a seguir, resolve o problema proposto no jogo.

% figura presente no enunciado

Calcule a razão $\dfrac{PS}{PQ}$.

---
tipo: Discursiva
dificuldade: Difícil
disciplina: Matemática
assunto: Teorema de Pitágoras
tags: [quadrado, área, Pitágoras]
fonte: material didático
gabarito:
resposta:
---
\question
A figura abaixo dá indicação sobre uma das centenas de provas do teorema de Pitágoras. Esta consiste em calcular as áreas dos quadrados construídos sobre os lados do triângulo $ABC$ e comparar os resultados.

% figura presente no enunciado

Considerando a configuração acima, se $AC$ mede $3\sqrt{3}$ cm e $BC$ mede 3 cm, determine, em cm, a medida de $BI$.

---
tipo: Discursiva
dificuldade: Difícil
disciplina: Matemática
assunto: Lei dos Cossenos
tags: [triângulo, lados ímpares consecutivos, UNICAMP]
fonte: concurso
banca: UNICAMP
ano: 2000
gabarito:
resposta:
---
\question
Os lados de um triângulo têm, como medidas, números inteiros ímpares consecutivos cuja soma é 15.
\begin{parts}
\part Quais são esses números?
\part Calcule a medida do maior ângulo desse triângulo.
\part Sendo $\alpha$ e $\beta$ os outros dois ângulos do referido triângulo, com $\beta > \alpha$, mostre que $\operatorname{sen}^2\!\beta - \operatorname{sen}^2\!\alpha < \dfrac{1}{4}$.
\end{parts}
